\chapter{Operations in depth}
\label{ch:ops}

An \emph{operation} is a recipe: it says \emph{how} to burn the artwork attached
to it. Five burn operation types exist, plus three effects that modify them and a
family of utility operations that move or pause the machine. This chapter is the
reference-tutorial for all of them: what they do, which settings matter, and how
to attach artwork to them.

\section{The five burn operations}

\begin{center}\small
\begin{tabular}{@{}L{1.8cm}L{1.4cm}L{3.2cm}L{3.2cm}L{4.3cm}@{}}
\toprule
\textbf{Type} & \textbf{Colour} & \textbf{Accepts} & \textbf{Ignores} &
\textbf{Use it for} \\
\midrule
\textbf{Cut} & red & paths, rectangles, ellipses, lines, polylines (and effects)
& images & cutting through: high power, slow speed, a few passes \\
\textbf{Engrave} & blue & the same vector shapes & images & scoring, marking,
thin vector engraving; dense hatching \\
\textbf{Raster} & black & vectors \emph{and} text, converted to a raster image &
images & filling an area quickly, engraving shapes and lettering \\
\textbf{Image} & (none) & image elements only & everything else & engraving a
raster image, with the image's own dither and contrast settings \\
\textbf{Dots} & (none) & points only & everything else & drilling, perforating,
dot marking: the head dwells in place instead of moving \\
\bottomrule
\end{tabular}
\end{center}

\noindent The ``Ignores'' column matters more than it looks. Assigning a
photograph to a \mkseries{Cut} operation gets you nothing, and by default a
\mkseries{Dots} or \mkseries{Image} operation \emph{stops} classification, so
elements landing there are not also picked up by later operations. When in doubt,
watch the operation's icon: unusable children are greyed out or flagged.

\begin{notebox}[Raster versus Image --- the difference that confuses everyone]
\mkseries{Raster} takes \emph{vector} artwork and turns it into a raster at job
time, using the operation's own DPI and direction settings. It is how you fill an
engraved logo or a rectangle of text. \mkseries{Image} takes an \emph{existing
raster image} and burns it, using the dither, contrast and halftone settings
stored on the image element itself. If the thing on your scene is a
\file{.png} or \file{.jpg}, you want \mkseries{Image}; if it is a path or a text
element, you want \mkseries{Raster}.
\end{notebox}

\section{Creating an operation}

Right-click the \emph{Operations} branch in the tree and choose \menu{Append
operation}, then one of: \btn{Append Image}, \btn{Append Raster},
\btn{Append Engrave}, \btn{Append Cut}, \btn{Append new Hatch} (which creates an
Engrave operation carrying a Hatch effect) or \btn{Append Dots}.

Notice that a new operation arrives with defaults taken from your device --- the
device can define what ``cut'' or ``engrave'' means for that machine, so a GRBL
diode machine and a CO\textsubscript{2} cutter do not start from the same
numbers. Typical starting points are around 5\unit{mm/s} for a cut,
25--35\unit{mm/s} for an engrave and 150--500\unit{mm/s} for a raster. These are
starting points, not answers: Chapter~\ref{ch:tuning} is about finding your own.

Under the same right-click menu, \menu{Append special operation(s)} adds the
non-burning steps. These run in the order they appear in the operations list,
which lets you build a job with a preamble (home, coolant on, wait for you to
press the machine's button) and a postscript (coolant off, return home).

\begin{itemize}[leftmargin=1.4em,itemsep=0pt]
  \item Movement and time: \btn{Append Home}, \btn{Append Physical Home},
        \btn{Append Return to Origin}, \btn{Append Goto Location}.
  \item Operator interaction: \btn{Append Beep}, \btn{Append Interrupt},
        \btn{Append Wait}, \btn{Append Output}, \btn{Append Input}.
  \item Ancillaries: \btn{Append Coolant On}, \btn{Append Coolant Off}.
  \item Combinations that do the above in one step, such as
        \btn{Append Home/Beep/Interrupt} and
        \btn{Append Origin/Beep/Interrupt}.
\end{itemize}

\section{Attaching artwork to an operation}

There are four ways, and you will settle into one of them within an hour.

\begin{enumerate}[leftmargin=1.6em]
  \item \textbf{Drag and drop.} Select elements in the scene, drag them onto the
        operation in the tree, and drop. The operation gains \emph{references} to
        those elements.
  \item \textbf{Right-click assign.} Right-click the selected elements in the
        scene, choose \menu{Assign Operation}, and pick the operation by name. The
        same submenu offers \btn{Remove all assignments from operations},
        \btn{Exclusive assignment}, \btn{Inherit stroke and classify similar} and
        \btn{Inherit fill and classify similar}.
  \item \textbf{Operation squares.} Open the \panel{Operations} pane
        (\menu{Panes>Tools>Operations}). It shows up to 24 coloured squares, one
        per operation. Left-click assigns using the element's \emph{stroke}
        colour as the classification colour; right-click uses the \emph{fill}
        colour instead. Options in the pane let you choose what a click does
        (\emph{Leave} / \emph{-> OP} / \emph{-> Elem}) and whether assignments are
        \emph{Similar} (colour-matched) or \emph{Exclusive}.
  \item \textbf{The operation list panel.} The compact operations panel shows one
        row per operation with an \emph{A} (active) toggle, an \emph{S} (visible)
        toggle, a coolant indicator, and editable \emph{Power} and \emph{Speed}
        columns --- the fastest way to tune a job by eye. Its
        \emph{Classification} section chooses whether colours are left alone, the
        operation inherits the element's colour, or the element inherits the
        operation's colour.
\end{enumerate}

\begin{warnbox}[Elements that are not assigned anywhere are invisible to the job]
An operation only burns what it references. A shape sitting in the elements
branch that is referenced by no operation will simply not be cut --- and
MeerK40t will warn you: ``Some shapes are not assigned to any operation \dots{}
Queue anyway?''. Learn to check for it: \menu{Edit>Select all non-assigned}
highlights the orphans. The \panel{Operation Info} window
(\menu{Tools>Operation Info}) lists every operation, its item count and its
runtime, and shows an \emph{Unassigned} row marked with \texttt{!} for elements
that belong to nothing.
\end{warnbox}

\section{Colour classification}

Each operation has a colour, and MeerK40t can use colour to decide where artwork
belongs. This is how an imported SVG with a black outline and a red fill becomes
``cut the red things, engrave the black things'' without you selecting anything.

\begin{itemize}
  \item A \textbf{Cut} operation is red by default. Elements whose stroke or fill
        matches that red are classified into it.
  \item An \textbf{Engrave} operation is blue by default.
  \item A \textbf{Raster} operation is black by default.
\end{itemize}

\noindent Classification is configured in
\menu{Edit>Settings>Preferences}, tab \tabname{Classification}, which offers two
presets: \btn{Automatic} and \btn{Manual}. Automatic classification on load is
what people usually want for production work with colour-keyed SVG files;
turn it off when you want to place everything by hand.

\begin{tipbox}
If colours confuse you, make them honest: in the element's own properties set its
stroke to the operation colour you intend to use. Then the scene itself shows the
plan --- red lines will be cut, blue lines engraved --- and mistakes are visible
at a glance.
\end{tipbox}

\section{The settings of a burn operation}

Select an operation in the tree and its properties appear in the \panel{Properties}
window (\menu{Tools>Properties}, or double-click). The settings below are the ones
that actually change results.

\begin{center}\small
\begin{tabular}{@{}L{3.2cm}L{2.2cm}L{8.7cm}@{}}
\toprule \textbf{Label} & \textbf{Default} & \textbf{Meaning} \\
\midrule \btn{Speed} & 10\unit{mm/s} (cut), 35 (engrave), 150 (raster/image), 35 (dots) & Head travel speed while burning. Shown as \unit{mm/min} instead when the device uses that convention. \\
\btn{Power} & 1000 & Power on a 0--1000 scale, or as a percentage when the device displays percentages. \\
\btn{Frequency} & 20\unit{kHz} & Pulse frequency. Only shown for machines that expose it. \\
\btn{Passes} & off, 1 pass & Repeat the whole operation. Tick \emph{custom} to set a count. \\
\btn{Kerf compensation} & 0\unit{mm} & Offset every path by half the kerf so finished parts come out the intended size (Chapter~\ref{ch:tuning}). \\
\btn{Enable} & on & Include this operation when the job is built. \\
\btn{Stop} & off (on for Image/Dots) & Stop classification here: later operations do not also capture these elements. \\
\btn{Coolant} & 0 = no change & 0 no change, 1 turn on, 2 turn off, for machines with coolant or air-assist control. \\
\btn{Color} & red / blue / black & The classification colour. \\
\btn{Loops} & 1 & Loop the operation branch (used with effects). \\
\bottomrule
\end{tabular}
\end{center}

\noindent The raster and image operations add their own settings:

\begin{center}\small
\begin{tabular}{@{}L{3.2cm}L{2.2cm}L{8.8cm}@{}}
\toprule \textbf{Label} & \textbf{Default} & \textbf{Meaning} \\
\midrule \btn{DPI} (with \emph{overrule}) & 500\unit{dpi} & Scan-line density. Higher = finer and much slower. \\
\btn{Overscan} & 0.5\unit{mm} & Extra travel at each end of a scan line so the head is already at speed when the beam starts. \\
\btn{Direction} & Bottom-to-top & Scan direction: top-to-bottom, bottom-to-top, right-to-left, left-to-right, crosshatch, greedy horizontal/vertical, crossover, spiral or diagonal. \\
\btn{Directional Raster} & Bidirectional & Burn on the way back as well, doubling throughput; uncheck for a directional look (or when your machine cannot handle it accurately). \\
\btn{Use grayscale instead} & on & Rasterise a black-and-white image as grayscale --- the raster op's answer to ``Override black/white image''. \\
\btn{Consider laserdot} & off & Skip pixels already covered by the laser's spot; saves time at low DPI. \\
\btn{Start Preference} & both on & Which corner of the raster the head starts from. \\
\btn{Dwell Time (ms)} & 50\unit{ms} & (Dots) how long the head rests on each point. \\
\bottomrule
\end{tabular}
\end{center}

\noindent DPI and raster step are two views of the same number: the step is the
distance between scan lines, and DPI is how many of them fit in an inch. At
500\unit{dpi} the step is 0.0508\unit{mm} --- about twice the spot size of a
typical 40\unit{W} CO\textsubscript{2} machine. Pushing DPI far beyond the spot
size buys darkness, not detail, and costs time linearly.

\section{Effects: hatch, wobble and warp}

Effects are children of a Cut or Engrave operation. They modify how that
operation's paths are burned.

\subsection*{Hatch}

\menu{Add effect>Add hatch effect} fills a closed shape with parallel lines
instead of tracing its outline --- the classic way to engrave a solid area with a
\emph{vector} (not raster) operation.

\begin{center}\small
\begin{tabular}{@{}L{3.1cm}L{1.9cm}L{9.4cm}@{}}
\toprule \textbf{Label} & \textbf{Default} & \textbf{Meaning} \\
\midrule \btn{Hatch Distance} & 1\unit{mm} & Spacing between hatch lines: smaller is darker and slower. \\
\btn{Angle} & 0\unit{deg} & Direction of the lines. \\
\btn{Angle Delta} & 0\unit{deg} & Rotate the pattern on each loop: 90\unit{deg} with 2 loops gives a cross-hatch. \\
\btn{Loops} & 1 & Repeat the fill, rotating per angle delta. \\
\btn{Unidirectional} & off & Burn every line in the same direction instead of back-and-forth. \\
\btn{Include Outline} & off & Also trace the shape's outline. \\
\btn{Fill Style} & scanline & Scanline, Eulerian or spiral fill algorithms. \\
\btn{Algorithm} & auto & \emph{Auto-Select}, \emph{Scanbeam (Default)} or \emph{Direct Grid (Fast)}. Jumping to a specific algorithm is a speed fix for pathological shapes. \\
\bottomrule
\end{tabular}
\end{center}

\subsection*{Wobble}

\menu{Add effect>Add wobble effect} replaces a path with a decorative oscillation
--- a wave, a row of circles, a gear feel --- by moving the head side to side as
it advances. Used for leather stitching effects, cork, and lines that should look
scratched or hand-drawn.

\begin{center}\small
\begin{tabular}{@{}L{3.1cm}L{2.0cm}L{9.3cm}@{}}
\toprule \textbf{Label} & \textbf{Default} & \textbf{Meaning} \\
\midrule \btn{Wobble Radius} & 0.5--1.5\unit{mm} & Amplitude of the sideways movement. \\
\btn{Wobble Interval} & 0.1\unit{mm} & How far the head travels before the pattern repeats. \\
\btn{Wobble Speed} & 50 & How quickly the pattern revolves along the path. \\
\btn{Fill Style} & circle & Circle, left/right circle, sinewave, sawtooth, jigsaw, gear, slowtooth or one of three meanders. \\
\bottomrule
\end{tabular}
\end{center}

\subsection*{Warp}

\menu{Add effect>Add warp effect} distorts the operation's artwork to fit a
containing shape you edit with the finger tool --- useful for engraving onto
curved or irregular objects where a straight projection would not follow the
surface.

\begin{notebox}[Effects hide their children]
A wobble effect, in particular, is drawn as the wobbled result rather than as
the original path. If a shape looked right before you added the effect and
strange after, remove the effect and compare; \menu{View>Scene Appearance} and
the operation's \emph{S} (visible) toggle control what you see without changing
what burns.
\end{notebox}

\section{Estimating before you commit}

Two windows tell you what a job will cost in time:

\begin{itemize}
  \item \panel{Operation Info} (\menu{Tools>Operation Info}) lists each
        operation with its type, name, item count and runtime; press \btn{Get Time
        Estimates} to fill in the last column. Unassigned elements appear as their
        own row, marked \texttt{!}.
  \item \panel{Simulation} (\btn{Simulate} on the Laser tab, or \menu{Tools}
        \menu{Simulation}) replays the job on screen and reports
        \emph{Laser Time}, \emph{Travel Time}, \emph{Total Time},
        \emph{Extra Time}, and the corresponding distances.
        Chapter~\ref{ch:order} uses this window properly.
\end{itemize}

\begin{machinebox}{Burn scrap first, always}
For any new material or any setting you have not personally verified on this
machine, run the job on scrap with the lid closed and watch it. The estimate is
arithmetic done on your settings; the burn is physics. They agree most of the
time.
\end{machinebox}

\begin{tryit}{Build a two-operation test plaque}
\begin{enumerate}[label=\arabic*.,leftmargin=1.4em]
  \item In the console, make two shapes:
        \cmd{rect 10mm 10mm 40mm 40mm} then \cmd{circle 80mm 30mm 15mm}.
  \item Right-click \emph{Operations} in the tree, \menu{Append operation>Append
        Cut}. Note it appears red.
  \item Right-click \emph{Operations} again, \menu{Append operation>Append
        Engrave}. Note it appears blue. Now two operations exist and the first is
        ``exclusive'' by default --- only one of them will own the shapes.
  \item Select both shapes in the scene (\key{Ctrl+A}), open the
        \panel{Operations} pane, and click the blue \emph{Engrave} square.
  \item Select only the circle, and click the red \emph{Cut} square.
  \item Open \menu{Tools>Operation Info} and press \btn{Get Time Estimates}.
        You should see both operations with one item each.
  \item Select the rectangle's operation and change its \emph{Speed} and
        \emph{Power}. Watch the estimate change.
\end{enumerate}
\end{tryit}
